\subsection*{Relación con trabajos previos}

Luria y Simkin~\cite{LS} obtuvieron una cota inferior para el conteo de reinas mediante un
proceso voraz y absorción. Bowtell y Keevash~\cite{BK} estudiaron el problema toroidal y sus
consecuencias para el tablero ordinario. Aquí la construcción del predictor parte de otro
objeto, el máximo de una entropía sobre pesos de casilla.

La aproximación de Bethe y sus correcciones por lazos se han estudiado para el permanente de matrices
no negativas~\cite{Vontobel,WC} y, en general, para modelos discretos sobre grafos~\cite{CC,WF,YFW}.
La fórmula de Ihara--Bass~\cite{Bass}, en la forma de Watanabe y Fukumizu~\cite{WF}, relaciona la
corrección de un lazo con un determinante de tamaño igual al número de líneas, y es el origen
heurístico del pseudodeterminante de~\eqref{eq:Qpred}.
Para las semirreinas toroidales, Eberhard, Manners y Mrazović~\cite{EMM19} obtuvieron una asintótica
que incluye el factor constante, y después trataron las transversales de cuadrados latinos
cuasialeatorios~\cite{EMM}; en esos modelos no intervienen las dos familias de diagonales del tablero
ordinario. Liu, Liao y Wang~\cite{LLW} estudiaron el problema de las reinas como un gas reticular, con
simulaciones de Monte Carlo y redes tensoriales.
Este artículo usa los resultados de~\cite{Perfil} sobre el perfil de masa uno. En~\cite{Completabilidad}
se estudian, mediante una comparación de entropía relativa con un proceso de referencia y una doble
cuenta de reparaciones, las colocaciones legales elegidas uniformemente a las que falta un número fijo
de reinas: la fracción media de cruces vacíos (las intersecciones de las filas y las columnas vacías)
que admiten una reina tiende al inverso de la constante de normalización de ese perfil. Ese resultado
no determina el término constante del conteo exacto.

\subsection*{Organización}

Las Secciones~\ref{sec:preliminares} y~\ref{sec:predictor} pasan de las cuatro líneas de una
casilla a la fórmula del predictor. A partir de ahí, las dos preguntas del artículo se tratan por
separado. Para obtener la constante, la Sección~\ref{sec:perfil} compara el tablero finito con el
perfil continuo y la Sección~\ref{sec:constante} calcula los términos que sobreviven en el límite.
Para entender la consistencia, la Sección~\ref{sec:insercion} estudia una inserción y las
Secciones~\ref{sec:macro} y~\ref{sec:meso} controlan su efecto en los dos regímenes del teorema.

La Sección~\ref{sec:numerico} contrasta dos evaluaciones de la constante del predictor con los
valores del conteo exacto para $16\le n\le27$. La Sección~\ref{sec:discusion} distingue esos resultados
de la cuestión abierta sobre $Q(n)$. El Apéndice~\ref{sec:torres} ofrece un caso explícito
en el que se conservan sólo las restricciones de filas y columnas.
